\chapter{Discussion}
\label{ch:discussion}

This chapter interprets the results of Chapters~\ref{ch:verification}
and~\ref{ch:results}, addressing the accuracy of the numerical
implementation, the conditions under which the guiding-centre
approximation holds, and the physical content of the planetary
magnetosphere simulations.

\section{Validity of the Numerical Implementation}
\label{sec:disc_validity}

The verification hierarchy in Chapter~\ref{ch:verification} demonstrates
that the solver correctly implements the Lorentz force equation across a
range of field configurations. In all cases with $\mathbf{E} = \mathbf{0}$
and static $\mathbf{B}$, kinetic energy is conserved to better than
$10^{-6}$ relative error. The analytic overlays in the verification tests
are reproduced to high precision: the maximum radial deviation in circular
gyromotion is below $10^{-8}$, the parallel motion in helical flight is
exact to $10^{-9}$, and the drift speeds in the $\mathbf{E} \times
\mathbf{B}$ and gradient-$B$ tests match theory to within 3\,\%. The
dipole field passes both the on-axis magnitude scaling and field-direction
checks. The tolerances $\texttt{rtol} = 10^{-9}$,
$\texttt{atol} = 10^{-12}$ (equation~\eqref{eq:tolerances_num}) are
therefore appropriate for all simulations in this project.

\section{When Does Guiding-Centre Theory Apply?}
\label{sec:disc_gc}

The drift results in the verification tests show that guiding-centre
theory accurately describes the slow motion of the guiding centre when
$\rho / L_{\nabla B} \ll 1$. In the gradient-$B$ test the gradient is
intentionally weak ($\varepsilon = 0.05$), placing the particle firmly in
the adiabatic regime; the numerical drift agrees with the analytic
prediction to within 3\,\%.

The direct comparison between the full Lorentz orbit and the GC
integrator in Section~\ref{sec:gc_vs_full} confirms this further: with
$\rho / r_\mathrm{eq} \approx 0.006$ and $T_b / T_\Omega \approx 100$,
the GC trajectory tracks the gyro-averaged full orbit across all four
bounce cycles with no visible phase error in the bounce envelope.

Guiding-centre theory is expected to break down when the Larmor radius
becomes comparable to the field-variation scale. For the dipole
simulations in Chapter~\ref{ch:results}, the approximate conservation of
$\mu$ below 1\,\% relative variation confirms that the chosen parameters
sit comfortably in the adiabatic regime. If the particle energy were
increased or the particle moved to smaller $L$-shells where field
gradients are steeper, $\mu$ conservation would degrade and the full
orbit integration would be necessary even to describe the slow drift
motion.

\section{Guiding-Centre Approximation versus Full Orbit}
\label{sec:disc_gc_full}

The results of Section~\ref{sec:gc_vs_full} show that in the deep
adiabatic regime ($\rho / r_\mathrm{eq} \approx 0.006$), the GC
integrator reproduces the slow bounce dynamics of the full Lorentz orbit
faithfully across approximately 400 gyrations, at a fraction of the
computational cost.

The slow growth of the positional offset between the extracted GC and
the GC equations solution is a consequence of secular phase drift in
the non-symplectic RK45 integrator rather than breakdown of the adiabatic
approximation: the local agreement in $z(t)$ and $v_\parallel(t)$
remains excellent throughout, and the envelope of $\mu$ shows no secular
trend. A symplectic integrator such as the Boris scheme would suppress
this drift and would be the natural choice for very long integration runs.

\section{Physical Units and the SI Simulation}
\label{sec:disc_si}

The SI simulation of Section~\ref{sec:si_units} serves two purposes: it
verifies that the dimensional factors are restored correctly, and it
tests the solver under conditions where the timescale separation is
extreme. At $L = 3$ a 10\,keV electron has
$\rho / r_\mathrm{eq} \approx 10^{-5}$ and
$T_b / T_\Omega \approx 57{,}000$, placing it far deeper into the
adiabatic regime than any of the dimensionless test cases. The numerical
mirror latitude agrees with the analytic prediction to within 0.02\,\%,
confirming that the implementation handles SI-scale quantities without
loss of precision.

The relative energy drift of approximately 1.5\,\% over two bounce
periods is larger than in the dimensionless tests, where
$|\delta K|$ stays below $10^{-6}$. This is not a failure of the solver
but a consequence of integrating approximately 115{,}000 gyrations: the
local truncation error per step is comparable to the dimensionless cases,
but it accumulates over far more cycles. The same mechanism produces the
2\,\% relative variation in $\mu$. Both quantities grow secularly rather
than oscillating, which is the signature of phase drift in a
non-symplectic integrator. The physically meaningful diagnostic is the
mirror-point accuracy, which remains excellent throughout, confirming
that the adiabatic invariant is well conserved in the underlying physics
even when the numerical representation of phase drifts over many
gyrations.

For applications requiring long-term phase accuracy at physical scales,
a symplectic integrator would be the appropriate choice. The RK45 scheme
used here is sufficient for the few-bounce-period runs in this project.

\section{Tilted Dipole Fields}
\label{sec:disc_tilted}

The tilted dipole results of Section~\ref{sec:tilted} confirm that the
solver handles the generalised dipole field correctly and that the mirror
mechanism operates as expected even when the magnetic and geographic
coordinate systems are misaligned. The key observation is that the
asymmetric $z(t)$ bounce seen at $47^\circ$ and $59^\circ$ tilt is not
a physical asymmetry in the particle dynamics but a geometric one: in
magnetic coordinates the bounce remains symmetric, and the mirror
latitudes are unchanged from the aligned case. The asymmetry appears
only when the trajectory is projected onto the geographic $z$-axis,
because the magnetic equatorial plane is tilted with respect to it.

This distinction matters for interpreting observations of trapped
particles at planets with large magnetic tilt. At Neptune ($47^\circ$
tilt) and Uranus ($59^\circ$ tilt), a spacecraft measuring particle
fluxes in geographic coordinates would see an apparently asymmetric
pitch-angle distribution even though the underlying bounce dynamics are
symmetric about the magnetic equator. The simulation reproduces this
effect directly: the bounce period and $\mu$ conservation are
indistinguishable across the three tilt angles, while the geographic
$z(t)$ traces look qualitatively different.

\section{\texorpdfstring{Corotation and $\mathbf{E} \times \mathbf{B}$ Drift}{Corotation and E x B Drift}}
\label{sec:disc_corotation}

The corotation test in Section~\ref{sec:corotation} provides a clean
decomposition of the azimuthal drift into its magnetic and electric
components. Running two simulations with identical initial conditions
--- one with and one without the corotation electric field --- isolates
the $\mathbf{E} \times \mathbf{B}$ contribution directly. The
difference recovers the input rotation rate $\Omega = 0.02$ to within
96.6\,\%, which is a strong confirmation that the corotation field is
implemented correctly.

The 3.4\,\% shortfall is physical rather than numerical. The corotation
velocity $\mathbf{v}_\mathrm{rot} = \boldsymbol{\Omega} \times
\mathbf{r}$ has a small component along $\mathbf{B}$ at non-equatorial
latitudes, and since the $\mathbf{E} \times \mathbf{B}$ drift is the
perpendicular part of $\mathbf{v}_\mathrm{rot}$, the effective drift
speed is slightly less than $|\mathbf{v}_\mathrm{rot}|$. A particle
that stayed exactly in the equatorial plane would co-rotate at exactly
$\Omega$, but the bouncing particle samples non-equatorial field
geometries where this projection effect reduces the time-averaged
drift. This is a real physical effect that the simulation captures
correctly.

\section{The Rotating Tilted Dipole}
\label{sec:disc_rotating}

The rotating tilted dipole of Section~\ref{sec:rotating} is the most
complete simulation in this project. It combines a time-dependent
magnetic field, corotation electric field, and three-dimensional
particle dynamics into a single integration. The fact that the solver
handles this case without modification --- simply evaluating
$\mathbf{B}(\mathbf{r}, t)$ and
$\mathbf{E}(\mathbf{r}, t)$ at each timestep --- validates the
generality of the implementation.

Two features of the results deserve comment. First, the modulation of
the bounce amplitude in $z(t)$ (Figure~\ref{fig:rot_z}) is a direct
consequence of the rotating field geometry: as the tilted dipole sweeps
past the particle, the effective mirror latitude shifts, producing an
amplitude envelope at the rotation period. This modulation is absent
in the static tilted case (Section~\ref{sec:tilted}), confirming
that it arises from the time dependence of the field rather than the
tilt itself.

Second, the measured azimuthal drift rate $\Omega_\mathrm{meas}
\approx 0.026$ exceeds the rotation rate $\Omega = 0.02$, consistent
with the gradient and curvature drift contribution measured
independently in the aligned corotation test
(Section~\ref{sec:corotation}). The additional gradient and curvature
drifts add to the corotation drift, producing a total azimuthal drift
rate that slightly exceeds the planetary rotation rate. This is the
expected physical behaviour for trapped particles in a rotating
magnetosphere.

Together, these results demonstrate that the solver can handle the full
complexity of a rotating planetary magnetosphere, reproducing all three
adiabatic motions (gyration, bounce, drift) simultaneously in a
time-varying field geometry.
